% Generated draft: requires mathematical and source review.
\documentclass[11pt]{article}
\usepackage[margin=25mm]{geometry}
\usepackage{fontspec}
\setmainfont{Latin Modern Roman}
\usepackage{amsmath,amssymb,mathtools,unicode-math}
\setmathfont{Latin Modern Math}
\usepackage{enumitem,needspace,xurl,hyperref}
\setlength{\parindent}{0pt}
\setlength{\parskip}{5pt}
\setlength{\emergencystretch}{3em}
\newcommand{\N}{\mathbb{N}}
\newcommand{\Z}{\mathbb{Z}}
\newcommand{\Q}{\mathbb{Q}}
\newcommand{\R}{\mathbb{R}}
\newcommand{\C}{\mathbb{C}}
\newcommand{\F}{\mathbb{F}}
\newcommand{\A}{\mathbb{A}}
\newcommand{\PP}{\mathbb{P}}
\newcommand{\E}{\mathbb{E}}
\newcommand{\eps}{\varepsilon}
\newcommand{\dd}{\,\mathrm{d}}
\begin{document}
\section*{UnsolvedMath ID: 30000929}
\textbf{Source problem code:} catalogue entry
\subsection{Definitions and mathematical statement}
For the matrix size and coefficient ring used in the cited report, let \(R\) be the polynomial ring in the independent entries of an upper-triangular Cartan matrix. Let \(\Gamma\) act on \(R\) through the exceptional mutations and allowed sign changes of that setting. Its invariant ring is
\[
R^\Gamma=\{r\in R:\gamma(r)=r\text{ for every }\gamma\in\Gamma\}.
\]
Write \(A\) for the coefficient ring and \((F_i)_{i\in I}\) for the explicitly constructed invariant polynomials referred to in the source. The question is whether
\[
R^\Gamma=A[F_i:i\in I].
\]
Thus every invariant should be a polynomial, with coefficients in \(A\), in the proposed \(F_i\). The supplied assessment reports that the \(F_i\) are invariant and algebraically independent. With that reported independence, the remaining generation question is equivalently whether the natural map
\[
A[t_i:i\in I]\longrightarrow R^\Gamma,\qquad t_i\longmapsto F_i,
\]
is surjective, and hence is the proposed polynomial-algebra presentation.

The record does not reproduce the matrix size, coefficient ring or characteristic, diagonal normalization and independent entries, mutation and sign-change formulas, or the list and definitions of the \(F_i\). These data are necessary to identify the particular invariant ring. In particular, no fixed diagonal value, Lie-theoretic Cartan convention, alternative mutation group, or replacement generators are inserted here. The statement is an exact generation conjecture for the report's construction, rather than the already reported independence or invariance assertion.
\subsection{Short English statement}
Do the proposed independent invariant polynomials generate every invariant of the exceptional-mutation and sign-change action on upper-triangular Cartan-matrix entries? The action and generators must be recovered from the report.
\subsection{Sources}
\begin{itemize}
\item Local UnsolvedMath registry, permanent ID \(30000929\), source code catalogue entry, ``Invariant Rings of Exceptional-Mutation Actions.'' No individual proposer is supplied.
\item \emph{Representation Theory of Finite Dimensional Algebras} (2008), Oberwolfach Reports, DOI 10.4171/owr/2008/09, pp. 401--472: \url{https://doi.org/10.4171/owr/2008/09}. This is the sole supplied external reference.
\item Supplied literature triage, checked 2026-08-21: status ``partially solved.'' It reports invariance and algebraic independence proved in the report, but full invariant-ring generation stated only as a conjecture, with no located general proof. This dated assessment derives from supplied OWR curation; time-sensitive status and cited claims require independent verification.
\end{itemize}
\end{document}
